\documentclass[10pt,a4paper]{amsart}
\usepackage{amsmath,amssymb,amsthm,graphicx}
\usepackage{amsthm,amsmath, amsfonts, amssymb, comment}
\newtheorem{theorem}{Theorem}[section]
\newtheorem{lemma}[theorem]{Lemma}
\newtheorem{proposition}[theorem]{Proposition}
\newtheorem{corollary}[theorem]{Corollary}
\newtheorem{question}[theorem]{Question}
\newtheorem*{conjecture}{Conjecture}
\theoremstyle{definition}
\newtheorem{definition}{Definition}[section]
\newtheorem*{programme}{Program}
\theoremstyle{remark}
\newtheorem*{remark}{Remark}
\usepackage{enumerate}
\title{Coding countable linear orders in triangle equivalence of bounded distance value sets}
\author{Mahmood Etedadialiabadi, Su Gao, Feng Li, Ruiwen Li}
\date{arxiv:2411.02760v2; CC BY 4.0}
\begin{document}
\maketitle
\noindent\textit{Editorial scope.} One counted original theorem. E\textasciicircum{}infinity\_\{S\_infinity\} is the standard universal S\_infinity-orbit equivalence relation; completeness of countable-linear-order isomorphism is an allowed input. The full author construction and inverse definability argument are included, with explicit distance-set and triangle-equivalence definitions.\par\medskip
\noindent\textit{Allowed standard prerequisites (editorial).}\begin{itemize}
\item Standard Borel coding of countable linear orders and a Borel embedding into the rational interval (3,4). An order on a countable enumerated domain can be embedded by placing successive elements between the finitely many previous images. The choice at each stage uses finitely many order relations; this is the routine coding used by source903/924.
\item Universal S\_infinity orbit equivalence relation. E\textasciicircum{}infinity\_\{S\_infinity\} denotes an equivalence relation Borel-universal for isomorphism of countable structures. Isomorphism of countable linear orders is a standard complete example.. Only the Consequently sentence of the counted theorem; the triangle encoding is not supplied by this fact.
\end{itemize}\medskip
Let $X$ be a set and let $\mathcal{B}$ be a $\sigma$-algebra of subsets of $X$. $(X, \mathcal{B})$ is a {\em standard Borel space} if there is a Polish topology $\tau$ on $X$ such that $\mathcal{B}$ is the $\sigma$-algebra of Borel sets generated (via iterations of countable union, countable intersection, and complementation) by $\tau$. Any Polish space is a standard Borel space. If $X$ is a standard Borel space and $Y$ is a Borel subset of $X$, then $Y$ is a standard Borel space. A subset $A$ of a standard Borel space is {\em analytic} if there is a standard Borel space $Y$ and there is a Borel function $f\colon Y\to X$ with $A=f(Y)$. An equivalence relation $E$ on a standard Borel space $X$ is {\em Borel} (or {\em analytic}) if $E$ is a Borel (analytic) subset of $X^2$. 

For equivalence relations $E, F$ on standard Borel spaces $X, Y$, respectively, we say that $E$ is {\em Borel reducible} to $F$, denoted $E\leq_B F$, if there is a Borel function $f\colon X\to Y$ such that for all $x, x'\in X$, we have
$$ xEx'\iff f(x)Ff(x'). $$
We say that $E$ is {\em Borel bireducible} with $F$, denoted $E\sim_B F$, if both $E\leq_B F$ and $F\leq_B E$ hold. We say that $E$ is {\em strictly Borel reducible} to $F$, denoted $E<_B F$, if $E\leq_B F$ but $F\not\leq_B E$. The Borel reducibility preorder $\leq_B$ gives rise to a relative complexity hierarchy for equivalence relations on standard Borel spaces. Specifically, given an equivalence relation $E$, we attempt to determine its exact complexity by comparing it, in terms of Borel reducibility, with some benchmark equivalence relations previously studied  and understood in the theory. If there is a benchmark equivalence relation $F$ so that we can show $E\sim_B F$, then we consider the exact complexity of $E$ to be completely determined. Short of doing this, we can establish upper bounds $E\leq_B F$ or $E<_B F$ for $E$, or lower bounds $F\leq_B E$ or $F<_B E$ for $E$, as much as we can, which provide partial information on the complexity of $E$.

    \begin{definition}[\cite{EGLMM}]
        Let $\Delta$ be a nonempty set of positive real numbers. We call $\Delta$ a {\em distance value set} if for all $x, y\in \Delta$, 
        \begin{equation*}
           \min\{ x+y, \sup\Delta\} \in \Delta.
        \end{equation*}
    \end{definition}
Briefly speaking, a distance value set is exactly an initial segment of a subsemigroup of positive reals. Given a distance value set $\Delta$, a metric space $(X,d)$ is a {\em $\Delta$-metric space} if $d(x,y)\in\Delta$ for all $x\neq y\in X$. 

\begin{definition}[\cite{EGLMM}]\label{def:CDVSequiv} Let $\Delta$ and $\Lambda$ be distance value sets. 
\begin{enumerate}
\item[(i)] A triple $(x,y,z)$ of positive real numbers is a {\em $\Delta$-triangle} if $x,y, z\in \Delta$ and $|x-y|\leq z\leq x+y$.
\item[(ii)] We say that $\Delta$ and $\Lambda$ are {\em equivalent} if there is a bijection $f\colon \Delta\to \Lambda$ such that for any positive real numbers $x,y, z$, $(x,y,z)$ is a $\Delta$-triangle if and only if $(f(x), f(y), f(z))$ is a $\Lambda$-triangle.
\end{enumerate}
\end{definition}

\begin{definition}
        Let $\mathsf{CDVS}\subseteq \mathbb{R}^{\mathbb{N}}$ consist of all sequences $(d_i)_{i\in\mathbb{N}}$ of real numbers such that
\begin{enumerate}
\item[(a)] $d_i\geq 0$ for all $i$; the set $\{i\in\mathbb{N}\colon d_i=0\}$ is infinite; $d_i>0$ for some $i$;
\item[(b)] if $d_i, d_j>0$, then $d_i=d_j$ if and only if $i=j$;
\item[(c)] if $\sup_kd_k<\infty$, then $\sup_k d_k=d_{i}$ for some $i$;
\item[(d)] for any $i,j\geq 0$, if $d_i+d_j<\sup_kd_k$ then there is $\ell$ such that $d_{\ell}=d_i+d_j$.
\end{enumerate}
\end{definition}

Intuitively, $\mathsf{CDVS}$ consists of codes of all countable distance value sets, which are enumerated as sequences of real numbers without repetitions. We include zero entries in the codes to accommodate finite distance value sets. $\mathsf{CDVS}$ is a Borel subset of $\mathbb{R}^{\mathbb{N}}$, and hence it is a standard Borel space. We also use the following notation:
$$\begin{array}{lcl}
\mathsf{CDVS}_0&=&\left\{(d_i)_{i\in\mathbb{N}}\in \mathsf{CDVS}\colon \inf\{d_i>0\colon i\in\mathbb{N}\}=0 \right\}, \\
\mathsf{CDVS}_+&=&\left\{(d_i)_{i\in\mathbb{N}}\in \mathsf{CDVS}\setminus \mathsf{CDVS}_0\colon \sup_i d_i<\infty\right\},\\
\mathsf{CDVS}_\infty &=& \left\{(d_i)_{i\in\mathbb{N}}\in \mathsf{CDVS}\setminus \mathsf{CDVS}_0\colon \sup_i d_i=\infty\right\}. 
\end{array}$$
These sets form a partition of $\mathsf{CDVS}$ into Borel subsets, and in particular they are all standard Borel spaces. $\mathsf{CDVS}_0$ consists of codes of all countable distance value sets $\Delta$ with $\inf\Delta=0$. Similarly, $\mathsf{CDVS}_+$ codes countable distance value sets $\Delta$ with $\inf\Delta>0$ and $\sup\Delta<\infty$, and $\mathsf{CDVS}_\infty$ codes those $\Delta$ with $\inf\Delta>0$ and $\sup\Delta=\infty$. In our discussions below, we will not distinguish a countable distance value set $\Delta$ from a code of $\Delta$ as an element of $\mathsf{CDVS}$.


\begin{definition}\label{def:er} %Let $\sim$ be the equivalence relation on $\mathbb{R}^{\mathbb{N}}$ defined by
%$$ (c_i)_i\sim (d_i)_i\iff \exists g\in S_\infty\ \exists r\in (0,+\infty)\ \forall i\ d_{g(i)}=r\cdot c_i. $$
Let $\approx$ be the equivalence relation on $\mathbb{R}^{\mathbb{N}}$ defined by
$$\begin{array}{rcl} (c_i)_i\approx (d_i)_i&\iff& \exists g\in S_\infty\ [\forall i\ (c_i=0\longleftrightarrow d_{g(i)}=0) \mbox{ and } \\
& & \forall i, j, k\ (|c_j-c_k|\leq c_i\leq c_j+c_k\longleftrightarrow \\
& & |d_{g(j)}-d_{g(k)}|\leq d_{g(i)}\leq d_{g(j)}+d_{g(k)})].
\end{array}
$$
For any $A\subseteq \mathbb{R}^{\mathbb{N}}$, %let $\sim_A$ denote $\sim\upharpoonright (A\times A)$, and similarly, 
let $\approx_A$ denote $\approx\upharpoonright (A\times A)$.
\end{definition}

\begin{lemma}\label{lem:isomorphism} $\approx_{\mathsf{CDVS}}$ is Borel reducible to an isomorphism relation of countable structures. Consequently, $\approx_{\mathsf{CDVS}}$ is Borel reducible to $E^\infty_{S_\infty}$.
\end{lemma}

\begin{proof} To each $\Delta\in \mathsf{CDVS}$ we associate a countable structure $\mathcal{A}_\Delta=(\Delta, R^{\mathcal{A}})$, where $R^{\mathcal{A}}$ is a ternary relation defined as
$$ R^{\mathcal{A}}(d_1,d_2,d_3)\iff \mbox{$(d_1,d_2,d_3)$ is a $\Delta$-triangle.}$$
It is clear that for $\Delta, \Lambda\in \mathsf{CDVS}$, we have $\Delta\approx \Lambda$ if and only if $\mathcal{A}_\Delta\cong \mathcal{A}_{\Lambda}$. This gives a Borel reduction from $\approx_{\mathsf{CDVS}}$ to an isomorphism relation of countable structures.
\end{proof}

\begin{theorem}\label{thm:CDVS+} The isomorphism relation of all countable linear orders is Borel reducible to $\approx_{\mathsf{CDVS}_+}$. Consequently, $E^\infty_{S_\infty}$ is Borel reducible to $\approx_{\mathsf{CDVS}_+}$.
\end{theorem}

\begin{proof}
For a countable linear order $\mathcal{O}$, let $\psi(\mathcal{O})\subseteq (3,4)$  be an isomorphic copy of $\mathcal{O}$, and let 
$$\varphi(\mathcal{O})=\{2,4,6\}\cup\psi(\mathcal{O})\cup(2+\psi(\mathcal{O})). $$
Then $\varphi(\mathcal{O})\in \mathsf{CDVS}_+$. We verify that $\varphi$ is a Borel reduction from the isomorphism relation of all countable linear orders to $\approx_{\mathsf{CDVS}_+}$.

First suppose $\mathcal{O}_1$ and $\mathcal{O}_2$ are isomorphic countable linear orders. Then there is an order-preserving bijection $f$ between $\psi(\mathcal{O}_1)$ and $\psi(\mathcal{O}_2)$. For $x\in\varphi(\mathcal{O}_1)$, let 
$$F(x)=\left\{\begin{array}{ll} f(x),  & \mbox{ if $x\in\psi(\mathcal{O}_1)$,} \\
f(x-2)+2, & \mbox{ if $x\in 2+\psi(\mathcal{O}_1)$,} \\
x, & \mbox{ otherwise.}
\end{array}\right.
$$
It is straightforward to check that $F$ is a bijection between $\varphi(\mathcal{O}_1)$ and $\varphi(\mathcal{O}_2)$ witnessing their equivalence. Thus $\varphi(\mathcal{O}_1)\approx \varphi(\mathcal{O}_2)$.

Conversely, suppose $\varphi(\mathcal{O}_1)\approx \varphi(\mathcal{O}_2)$ and this equivalence is witnessed by a bijection $g$. Let $\Delta=\varphi(\mathcal{O}_1)$ and $\Lambda=\varphi(\mathcal{O}_2)$, and let $\mathcal{A}_\Delta=(\Delta, R^\Delta)$ and $\mathcal{A}_\Lambda=(\Lambda, R^\Lambda)$ be the countable structures defined in the proof of Lemma~\ref{lem:isomorphism}.  Then $g$ is an isomorphism witnessing $\mathcal{A}_\Delta\cong \mathcal{A}_\Lambda$. 

Now note that $2$ is definable in both $\mathcal{A}_\Delta$ and $\mathcal{A}_\Lambda$ by the formula $\exists y\neg R(x, x, y)$. Thus $g(2)=2$. 
Moreover, $\psi(\mathcal{O}_1)$ is definable in $\mathcal{A}_\Delta$ by the formula 
$$ x\neq 2\wedge R(2, 2, x) \wedge \exists y\neg R(2,x,y). $$
Similarly, $\psi(\mathcal{O}_2)$ is definable in $\mathcal{A}_\Lambda$ by the same formula. Thus $g$ sends $\psi(\mathcal{O}_1)$ to $\psi(\mathcal{O}_2)$. Finally, on $\psi(\mathcal{O}_1)$, the $x<y$ relation is definable by the formula
$$ \exists z\, [\neg R(2, x, z) \wedge R(2,y,z)]. $$
Similarly for $\psi(\mathcal{O}_2)$. Hence $g$ preserves the order $<$ as a bijection between $\psi(\mathcal{O}_1)$ and $\psi(\mathcal{O}_2)$. It follows that the linear orders $\mathcal{O}_1$ and $\mathcal{O}_2$ are isomorphic.

It is routine to check that $\varphi$ is Borel.
\end{proof}

\begin{thebibliography}{99}
\bibitem{EGLMM}
M. Etedadialiabadi, S. Gao, F. Le Ma\^{\i}tre, and J. Melleray,
\textit{Dense locally finite subgroups of automorphism groups of ultraextensive spaces},
Adv. Math. {\bf 391} (2021), 107966.


\bibitem{GaoBook}
S. Gao,
Invariant Descriptive Set Theory,
Monographs and Textbooks in Pure and Applied Mathematics, vol. 293,
CRC Press, 2009.

\end{thebibliography}
\end{document}
